
\section{Актуальность темы исследования}

Современные мобильные роботы --- беспилотные летательные аппараты (БПЛА), надводные и
наземные платформы --- широко применяются для картографирования, геодезии и инспекции
протяжённых промышленных объектов: трубопроводов, высоковольтных линий электропередачи,
мостовых конструкций. Базовым средством навигации служат глобальные навигационные
спутниковые системы (ГНСС) в режиме. Однако металлические
конструкции промышленных объектов создают зоны экранирования и многолучёвости сигнала,
при которых кратковременные потери ГНСС продолжительностью 3--240~с являются штатной
эксплуатационной ситуацией.

В отсутствие ГНСС позиционирование переходит на бортовую инерциальную навигационную
систему (БИНС) на микроэлектромеханических датчиках (МЭМС). При двойном интегрировании
шума акселерометра ошибка позиции нарастает нелинейно: на интервале 60~с погрешность
достигает нескольких метров, при 240~с --- сотен метров. Восстановление утраченного участка
траектории в режиме постобработки позволяет обеспечить непрерывность данных без
повторного сканирования.

Дополнительную трудность создаёт смена динамического режима: при компенсации порывов
ветра квадрокоптер производит интенсивные асимметричные изменения тяги отдельных роторов,
которые не отражаются в показаниях акселерометра. Единственным датчиком, непосредственно
фиксирующим эти управляющие воздействия, является телеметрия оборотов несущих винтов.
Традиционные методы (фильтры Калмана) применимы лишь до 3--5~с, а Transformer имеет
квадратичную сложность~$O(N^2)$, неприемлемую для бортовых систем малых роботов с
ограниченными вычислительными ресурсами.

Таким образом, разработка нейросетевого метода восстановления траектории, сочетающего
линейную вычислительную сложность, физическую согласованность и новый информационный
канал --- телеметрию оборотов несущих винтов, --- представляет актуальную
научно-техническую задачу в области робототехники.

\section{Степень разработанности темы исследования}

\subsection{Классические методы.} 
Расширенный фильтр Калмана (РФК), фильтр без запаха (ФБЗ)
и сглаживатель Рауха--Тунга--Штрибеля (СРТШ) обеспечивают высокую точность при потерях
до 3--5~с. На более длинных интервалах ошибка позиции нарастает нелинейно вследствие
двойного интегрирования шума акселерометра (О.\,А.~Степанов, В.\,Я.~Распопов
и др.). Ни один из классических методов не использует телеметрию актуаторов в качестве
навигационного источника.

\subsection{Нейросетевые методы.} 
Двунаправленная LSTM \cite{Guang2021} и Transformer Encoder \cite{Guyard2025} превосходят
фильтры Калмана по точности, однако LSTM подвержена затуханию градиентов, а Transformer
имеет квадратичную сложность~$O(N^2)$, что исключает применение на бортовых системах с
нейроускорителями\cite{Chen2024}.

\subsection{Физически информированные нейронные сети PINN}
Физически информированные нейронные сети \cite{raissi2019physics}
встраивают уравнения динамики объекта в функционал потерь.
В задачах навигации БПЛА исследованы И.\,А.~Бизяевым, Б.\,Г.~Ильясовым и соавт.;
показана их перспективность при ограниченном объёме данных. Работы по использованию
телеметрии оборотов несущих винтов в качестве навигационного источника (G.~Lu,
J.~Cui и соавт., J.~Svacha и соавт.) подтверждают принципиальную важность этого канала.

\textbf{Архитектуры SSM.} Нейронная сеть на основе последовательностей с
селективными пространствами состояний (НС~ПСПС; базовая работа Mamba: A.~Gu, T.~Dao)
обеспечивает линейную сложность~$O(N)$, применяется в задачах управления роботами, однако
ни в одной известной работе не исследовалась для постобработки траектории при потерях ГНСС
с физически информированным функционалом потерь на базе телеметрии оборотов.

\section{Цель и задачи работы}

\textbf{Цель}: разработка и исследование физически информированной нейросетевой системы
восстановления траектории мобильного робота по данным БИНС и телеметрии оборотов несущих
винтов при кратковременных потерях сигнала ГНСС.

\textbf{Задачи:}
\begin{enumerate}
	\item Проанализировать классические и нейросетевые методы восстановления траектории,
	выявить их ограничения при потерях 3--240~с.
	\item Разработать эффективную модель НС~ПСПС с линейной сложностью~$O(N)$.
	\item Разработать физически информированную модель ФИ-НС~ПСПС, использующую
	телеметрию оборотов как новый информационный канал и как основу динамического
	функционала потерь на базе закона тяги~$F{=}k{\cdot}n^2$.
	\item Создать синтетический датасет Gazebo/PX4 с синхронизированными данными БИНС,
	ГНСС и оборотов двигателей квадрокоптера.
	\item Провести комплексные экспериментальные исследования на синтетическом
	и реальном датасете AI-IO.
	\item Провести применимость разработанного метода на платформах с 
\end{enumerate}

\textbf{Объект исследования} --- система навигации мобильного робота (многороторного
беспилотного летательного аппарата), функционирующая в условиях кратковременных потерь
сигнала глобальной навигационной спутниковой системы продолжительностью от~3 до~240~секунд.

\textbf{Предмет исследования} --- методы и алгоритмы восстановления траектории мобильного
робота в режиме постобработки на основе данных бесплатформенной инерциальной навигационной
системы и телеметрии оборотов несущих винтов, включая физически информированные
нейросетевые модели с динамическим функционалом потерь.


\section{Научная новизна}

\textbf{1.} Впервые для задачи восстановления траектории мобильного робота при потерях сигнала ГНСС предложена и исследована нейросетевая архитектура отличающаяся использованием модели последовательностей с селективными пространствами состояний (ПСПС) обеспечивающая линейную сложность. Что допускает исполнение предложенной модели в реальном времени на бортовых нейроускорителях с ограниченными ресурсами. 

\textbf{2.} Разработан метод физически информированного обучения нейросетевой модели с функцией потерь физической невязки на основе оборотов несущих винтов, отличающийся от известных физически информированных нейронных сетей для БПЛА, где физические ограничения обычно формулируются на основе кинематических соотношений или только данных инерциальных датчиков. Данный метод позволяет получать более сильный и физически обоснованный обучающий сигнал именно в режимах интенсивного дифференциального управления тягой.

\textbf{3.} Разработан и обоснован подход к использованию телеметрии оборотов несущих винтов в качестве самостоятельного навигационного источника, независимого от бесплатформенной инерциальной навигационной системы, отличающиеся от существующих работ, в которых обороты двигателей применяются только как дополнительный признак для улучшения оценки скорости. Данный подход демонстрирует что дифференциальные изменения тяги роторов принципиально ненаблюдаемы акселерометром и гироскопом, а так же демонстрирует информационнцю состоятельность оборотов несущих винтов.

\textbf{4.} Предложен принцип режимозависимой физической регуляризации, при котором функционал потерь автоматически усиливает вклад физических невязок при переходе платформы в динамический режим с повышенной активностью приводов, отличающиеся от  классических физически информированных нейронных сетей с фиксированными весовыми коэффициентами физических членов потерь. Предложенный принцип обеспечивает адаптивное усиление обучающего сигнала именно в тех режимах, где данные инерциальных датчиков наименее информативны. Это позволяет достигать высокой точности при дообучении модели на минимальном объёме калибровочных данных нового динамического режима.

\textbf{4.} \textbf{TODO:} Здесь должен быть пункт про исследование на разных роботах

\section{Теоретическая и практическая значимость исследования}

\textbf{Теоретическая значимость} исследования заключается в следующем. Предложен
принцип использования телеметрии актуаторов мобильного робота как нового самостоятельного
навигационного источника, независимого от данных инерциальной системы, что расширяет класс
информационных каналов, применимых в нейросетевых системах навигации без ГНСС. Разработан
и теоретически обоснован метод режимозависимой физической регуляризации, обеспечивающий
автоматическое усиление обучающего сигнала при смене динамического режима платформы.
Предложенная архитектура ФИ-НС~ПСПС расширяет класс физически информированных нейронных
сетей на архитектуры с линейной вычислительной сложностью типа SSM, что открывает новое
направление совместного применения физических ограничений и последовательностных моделей
в задачах мобильной робототехники. Разработанный синтетический датасет с синхронизированными
данными БИНС, ГНСС и RPM обеспечивает воспроизводимую исследовательскую основу для
задачи восстановления траектории при потерях ГНСС.

\textbf{Практическая значимость} исследования состоит в следующем. Разработанная
программа ФИ-НС~ПСПС (свидетельство о государственной регистрации №~2026618441 от
25~марта~2026~г.) пригодна к непосредственному применению в системах постобработки данных
беспилотных инспекционных комплексов. После квантования модель занимает 28~КБ; расчётное
время инференса на нейроускорителе Neural-ART (STM32N6, 600~GOPS при INT8) не превышает
1~мс на сегмент, что обеспечивает работу в режиме реального времени при частоте 100~Гц.
На реальном датасете AI-IO при потере ГНСС 60~с достигается СКО~$2{,}88{\pm}0{,}22$~м.
Предложенный операционный протокол --- один двухминутный калибровочный полёт перед
миссией --- обеспечивает адаптацию модели к новому динамическому режиму непосредственно
в полевых условиях с СКО~$9{,}4{\pm}1{,}8$~м при потере 240~с, что в 1,8 раза лучше
базовой модели на полном датасете. Результаты применимы для беспилотных комплексов,
выполняющих инспекцию линейной инфраструктуры (трубопроводы, ЛЭП), где кратковременные
потери ГНСС являются штатной эксплуатационной ситуацией.

\section{Методология и методы исследования}

Теоретической основой служат три физических закона:

\begin{enumerate}
	\item \textbf{Закон тяги несущего винта:} тяга~$i$-го ротора $F_i = k n_i^2$,
	суммарная вертикальная тяга $F_z = k\textstyle\sum_{i=1}^4 n_i^2$ уравновешивает вес~$mg$.
	\item \textbf{Второй закон Ньютона} для поступательного движения:
	$m\ddot{\mathbf{p}}_I = m\mathbf{g} + R_{IB}\mathbf{F}_B + \mathbf{F}_{a,I}$,
	где~$\mathbf{F}_{a,I} = {-}b_\text{drag}m\mathbf{v}$.
	\item \textbf{Уравнения управляющих моментов:} $\tau_x = l(F_0{-}F_2)$,
	$\tau_y = l(F_1{-}F_3)$, $\tau_z = b(F_0{-}F_1{+}F_2{-}F_3)$,
	где~$l$~--- плечо, $b$~--- аэродинамическая константа момента.
\end{enumerate}

Реализация~--- Python/PyTorch. Симуляция квадрокоптера Holybro X500~--- Gazebo/PX4.
Обучение нейронных сетей выполнялось на GPU NVIDIA GeForce GTX 1650. Верификация~--- 10 независимых
тестовых сегментов по метрикам СКО и АТО.

\section{Положения, выносимые на защиту}

\textbf{1.} НС~ПСПС обеспечивает линейную сложность~$O(N)$ и превосходит Transformer по
скорости восстановления траектории в 2,6--30~раз при длительностях потери 3--120~с при
сопоставимой точности.

\textbf{2.} Включение нового информационного канала (телеметрия оборотов~$n_i$) и
физической модели динамики в функционал потерь обеспечивает СКО~$14{,}3{\pm}2{,}3$~м
при дообучении на 10\% данных --- ниже, чем у базовой модели на 100\%.

\textbf{3.} ФИ-НС~ПСПС превосходит базовые нейросетевые методы и метод РФК по
совокупности критериев «точность -- скорость -- ресурсоёмкость» при длительностях потери
ГНСС от~30 до~240~с на синтетическом датасете и при всех исследованных длительностях
(3--60~с) на реальном датасете AI-IO.

\section{Степень достоверности полученных результатов}

Достоверность полученных результатов обеспечена корректным применением методов
математической статистики, теории нейронных сетей и уравнений динамики квадрокоптера.
Метрики СКО и АТО вычислялись на независимых тестовых сегментах, составляющих 20\%
датасета и не участвовавших в обучении; в таблицах результатов приводятся значения
«среднее~$\pm$~стандартное отклонение» по 10 непересекающимся сегментам.
Физическая адекватность предложенной модели подтверждена верификацией на открытом
датасете AI-IO лаборатории SJTU-ViSYS, не участвовавшем в разработке метода; результаты
РФК для той же платформы использованы как внешний базовый показатель.
Скоростные характеристики НС~ПСПС подтверждены сопоставлением с Transformer Encoder
и двунаправленной LSTM на адаптированном датасете EuRoC~MAV.
Разработанная программа зарегистрирована в Федеральной службе по интеллектуальной
собственности (свидетельство №~2026618441 от 25~марта~2026~г.), что обеспечивает
независимую верификацию факта создания программного продукта.


\section{Личный вклад автора}

Соискателем самостоятельно сформулированы цель и задачи исследования. Разработаны
архитектуры моделей НС~ПСПС и ФИ-НС~ПСПС, включая физический функционал потерь
на основе телеметрии оборотов несущих винтов. Создан синтетический датасет на базе
симулятора Gazebo/PX4, содержащий синхронизированные данные БИНС, ГНСС и оборотов
четырёх двигателей. Проведены все экспериментальные исследования, выполнен анализ
и обобщение результатов. Разработана и зарегистрирована программа для ЭВМ. В совместных
публикациях автор принимал непосредственное участие в постановке задачи, разработке
метода, проведении экспериментов и обработке результатов. Все основные научные
положения, компьютерные модели, практические решения и теоретические выводы
сформулированы и выполнены автором самостоятельно.

\section{Соответствие диссертации паспорту научной специальности}

Диссертационная работа выполнена в соответствии с паспортом специальности 2.5.4
«Роботы, мехатроника и робототехнические системы» по пунктам:

\begin{itemize}
	\item п.~5. \textit{Методы, алгоритмы, программные и аппаратные средства управления
		роботами, робототехническими и мехатронными системами, включая адаптивное,
		оптимальное, распределённое, интеллектуальное и супервизорное управление.}
	
	\item п.~6. \textit{Математическое и программное обеспечение, компьютерные методы
		и средства обработки информации в реальном времени в роботах, робототехнических
		и мехатронных системах.}
\end{itemize}

\section{Структура и объём работы}

Диссертация состоит из введения, четырех глав, заключения, списка использованных источников
и двух приложений. Содержание работы изложено на~[N]~страницах, включает~[N]~рисунков
и~[N]~таблиц. Библиографический список содержит 26~наименований.


\section{Основные публикации по теме диссертации}
% =========================================================

\textbf{Публикации в изданиях, входящих в перечень ВАК:}

\begin{enumerate}
	\item Сидоренко~Д.\,Д., Майстро~А.\,С., Булдаков~П.\,Ю.,
	Сыромятников~А.\,Д., Тимофеев~А.\,Н.
	Сокращение вычислительных ресурсов в задаче восстановления траектории
	мобильных роботов при потере сигнала ГНСС
	// Наука и бизнес: пути развития. --- 2025.
	
	\item Сидоренко~Д.\,Д., Тимофеев~А.\,Н., Романов~П.\,И.
	Физически информированная нейронная сеть для восстановления траектории
	квадрокоптера в условиях потери сигнала ГНСС
	// (в печати).
\end{enumerate}

\textbf{Свидетельство о государственной регистрации программы для ЭВМ:}

\begin{enumerate}
	\setcounter{enumi}{2}
	\item Сидоренко~Д.\,Д. ФИ-НС~ПСПС~--- программа восстановления траектории
	квадрокоптера. Свидетельство о государственной регистрации программы для ЭВМ
	№~2026618441 от 25~марта~2026~г.
\end{enumerate}